\uuid{dPvd}
\niveau{PCSI}
\module{Analyse}
\chapitre{Éléments de logique et raisonnement}
\sousChapitre{Ensembles}
%!TeX root=../../../encours.nouveau.tex

\duree{5}
\difficulte{1}
\auteur{Antoine Crouzet}
\datecreate{2024-12-01}
\titre{Ensembles}
\contenu{
\question{Soient les ensembles suivants : \[A = \{1,\,2,\,3,\,4,\,5,\,6,\,7\} ,\quad  B = \{1,\,3,\,5,\,7\} , \quad C = \{2,\,4,\,6\} , \qeq D = \{3,\,6\}\]

Déterminer $B\cap D, C\cap D, B\cup C, B \cup D$. Déterminer les complémentaires dans $A$ de $B, C$ et $D$.}
\reponse{
Exercice rapide pour bien maitriser les différents concepts. On obtient :
\[B\cap D = \{3\},\quad C\cap D = \{6\},\quad B\cup C=\{1,\,2,\,3,\,4,\,5,\,6,\,7\}=A,\quad B\cup D = \{1,\,3,\,5,\,6,\,7\}\]
Enfin, dans $A$, on a :
\[\overline{B}=\{2,\,4,\,6\}=C,\quad \overline{C}=\{1,\,3,\,5,\,7\}=B,\quad \overline{D}=\{1,\,2,\,4,\,5,\,7\}\]
}
}
